\begin{center}
\begin{tikzpicture}[>=stealth]
 \draw[thick] (0,0) ellipse (4.85 and 4.14);
 \fill[orange!80] (2.525,0) circle (0.3);
 \node at (2.525,-0.6) {Sonne};
 \fill (4.85,0) circle (0.1);
 \node[anchor=west] at (4.95,0.3) {Perihel};
 \fill (-4.85,0) circle (0.1);
 \node[anchor=east] at (-4.95,0.3) {Aphel};
 \draw[|<->|] (2.525,-1.1) -- (4.85,-1.1) node[midway, below] {$r_\mathrm{P}$};
 \draw[|<->|] (-4.85,-1.1) -- (2.525,-1.1) node[midway, below] {$r_\mathrm{A}$};
 \draw[->, very thick] (4.85,0) -- (4.85,1.5) node[right] {$v_\mathrm{P}$};
\end{tikzpicture}
\end{center}

\begin{abcliste}
\abc Mit dem Nullpunkt an der Sonnenoberfläche ($R_\odot = \RSoO$, $M_\odot = \MSoO$, FoTaBe S.~211) ist die potentielle Energie im Abstand $r$
\begin{align}
 E_\mathrm{p}(r) &= -G\,\frac{M_\odot\,m}{r}-\left(-G\,\frac{M_\odot\,m}{R_\odot}\right) = G\,M_\odot\,m\left(\frac{1}{R_\odot}-\frac{1}{r}\right)
\end{align}
Im Perihel und im Aphel:
\begin{align}
 E_\mathrm{P} &= \EPeF \\
   &= \ncG\cdot\MSo\cdot\mHe\cdot\left(\frac{1}{\RSo}-\frac{1}{\rPe}\right) \\
   &= \EPe \approx \result{\EPeP{12}{3}} \\
 E_\mathrm{A} &= \EApF \\
   &= \ncG\cdot\MSo\cdot\mHe\cdot\left(\frac{1}{\RSo}-\frac{1}{\rAp}\right) \\
   &= \EAp \approx \result{\EApP{12}{3}}
\end{align}
Die beiden Werte unterscheiden sich nur um etwa ein Prozent, weil beide Punkte im Vergleich zum Sonnenradius sehr weit weg sind.

\abc Die Sonde muss mit ihrer kinetischen Energie im Perihel mindestens die Differenz der potentiellen Energien bis zum Aphel aufbringen können (im Aphel hätte sie dann gerade die Geschwindigkeit null):
\begin{align}
 \frac{1}{2}\,m\,v_\mathrm{P}^2 &\geq E_\mathrm{A}-E_\mathrm{P} = G\,M_\odot\,m\left(\frac{1}{r_\mathrm{P}}-\frac{1}{r_\mathrm{A}}\right)
\end{align}
Damit ist
\begin{align}
 v_\mathrm{P} &\geq \vPmF \\
   &= \sqrt{2\cdot\ncG\cdot\MSo\cdot\left(\frac{1}{\rPe}-\frac{1}{\rAp}\right)} \\
   &= \vPm \approx \result{\vPmP{3}{2}}
\end{align}
Tatsächlich fliegt die Sonde auch im Aphel noch. Mit dem zweiten Keplerschen Gesetz, $v_\mathrm{P}\,r_\mathrm{P} = v_\mathrm{A}\,r_\mathrm{A}$, und der Energieerhaltung erhält man die wirkliche Perihelgeschwindigkeit auf dieser Bahn:
\begin{align}
 v_\mathrm{P} &= \vPkF = \vPk \approx \vPkP{3}{2}
\end{align}
\end{abcliste}
